% GENERATED FILE. Edit canonical lessons, then run npm run generate:books.
% canonical-source-hash: fnv1a64:cd77dd2bb85617bf
% canonical-lessons: ML-C469-natappata, ML-C469-valiyoram, ML-C469-cuntupalaka, ML-C469-valimaruka, ML-C469-muriccukatakkuka

\chapter{Walking in Town}
\label{ch:ml-walking-in-town}
\hlchaptermodality{469}

\begin{chapteropening}
\textbf{By the end of this chapter:} \emph{I can say a footpath, the roadside, a signpost, to make way and to cross over.}

Complete the last lesson of chapter 469: I can say a footpath, the roadside, a signpost, to make way and to cross over.
\end{chapteropening}

\section[\texorpdfstring{\ml{നടപ്പാത} (naṭappāta)}{നടപ്പാത (naṭappāta)}]{\ml{നടപ്പാത} (naṭappāta) — a footpath; a pavement}

\label{lesson:ML-C469-natappata}



\begin{quote}
Before the new one: say the Malayalam for a flyover, then the Malayalam for to pull over.
\end{quote}

\subsection*{You'll want to know: \ml{നടപ്പാത}}
\textbf{\ml{നടപ്പാത}} — \emph{naṭappāta} — \textquotedblleft{}a footpath; a pavement\textquotedblright{}.

\begin{quote}
\ml{കുട്ടികൾ} \ml{നടപ്പാതയിലൂടെ} \ml{പള്ളിക്കൂടത്തിലേക്ക്} \ml{നടന്നു}.
\end{quote}

\begin{quote}
\emph{kuṭṭikaḷ naṭappātayilūṭe paḷḷikkūṭattilēkkŭ naṭannu.}
\end{quote}

\begin{quote}
The children walked to school along the footpath.
\end{quote}

Built from \ml{നട}, walking, the stem of \ml{നടക്കുക}, and \ml{പാത}, a road.

\begin{cousinweb}
\textbf{\ml{നടപ്പാത}} (\emph{naṭappāta}) is a Malayalam compound built from \ml{നട}, a walk, from the native root of \ml{നടക്കുക} (DEDR 3582), and \ml{പാത}, a road.
\end{cousinweb}

\begin{grammarlens}[title={What you've built}]
One more word for getting around.
\end{grammarlens}

\subsection*{Your turn}
\begin{itemize}
\raggedright
  \item \emph{Say it:} \emph{naṭappāta}
  \item \emph{Say it:} \emph{naṭappāta}, once more
  \item \emph{Say it:} \emph{kuṭṭikaḷ naṭappātayilūṭe paḷḷikkūṭattilēkkŭ naṭannu.} — The children walked to school along the footpath.
  \item \emph{Recall:} say the Malayalam for a flyover, then the Malayalam for to pull over, then say \emph{naṭappāta} again
\end{itemize}

\begin{tcolorbox}[breakable,colback=teal!4,colframe=teal!35!black,title={Before you move on}]
What does \ml{നടപ്പാത} mean? (\textquotedblleft{}A footpath; a pavement\textquotedblright{}.) Say it once more.
\end{tcolorbox}
\section[\texorpdfstring{\ml{വഴിയോരം} (vaḻiyōraṁ)}{വഴിയോരം (vaḻiyōraṁ)}]{\ml{വഴിയോരം} (vaḻiyōraṁ) — the roadside}

\label{lesson:ML-C469-valiyoram}



\begin{quote}
Before the new one: say the Malayalam for to pull over, then the Malayalam for a footpath.
\end{quote}

\subsection*{You'll want to know: \ml{വഴിയോരം}}
\textbf{\ml{വഴിയോരം}} — \emph{vaḻiyōraṁ} — \textquotedblleft{}the roadside\textquotedblright{}.

\begin{quote}
\ml{വഴിയോരത്ത്} \ml{ഒരു} \ml{ചെറിയ} \ml{കട} \ml{ഉണ്ട്}.
\end{quote}

\begin{quote}
\emph{vaḻiyōrattŭ oru ceṟiya kaṭa uṇṭŭ.}
\end{quote}

\begin{quote}
There is a small shop by the roadside.
\end{quote}

\ml{ഓരം} means an edge. A -\ml{ം} noun takes -\ml{ത്ത്} for at: \ml{വഴിയോരത്ത്}, at the roadside.

\begin{cousinweb}
\textbf{\ml{വഴിയോരം}} (\emph{vaḻiyōraṁ}) is a native Dravidian compound of \ml{വഴി}, a way (DEDR 5297), and \ml{ഓരം}, an edge or margin (DEDR 1060).
\end{cousinweb}

\begin{grammarlens}[title={What you've built}]
One more word for getting around.
\end{grammarlens}

\subsection*{Your turn}
\begin{itemize}
\raggedright
  \item \emph{Say it:} \emph{vaḻiyōraṁ}
  \item \emph{Say it:} \emph{vaḻiyōraṁ}, once more
  \item \emph{Say it:} \emph{vaḻiyōrattŭ oru ceṟiya kaṭa uṇṭŭ.} — There is a small shop by the roadside.
  \item \emph{Recall:} say the Malayalam for to pull over, then the Malayalam for a footpath, then say \emph{vaḻiyōraṁ} again
\end{itemize}

\begin{tcolorbox}[breakable,colback=teal!4,colframe=teal!35!black,title={Before you move on}]
What does \ml{വഴിയോരം} mean? (\textquotedblleft{}The roadside\textquotedblright{}.) Say it once more.
\end{tcolorbox}
\section[\texorpdfstring{\ml{ചൂണ്ടുപലക} (cūṇṭupalaka)}{ചൂണ്ടുപലക (cūṇṭupalaka)}]{\ml{ചൂണ്ടുപലക} (cūṇṭupalaka) — a signpost}

\label{lesson:ML-C469-cuntupalaka}



\begin{quote}
Before the new one: say the Malayalam for a footpath, then the Malayalam for the roadside.
\end{quote}

\subsection*{You'll want to know: \ml{ചൂണ്ടുപലക}}
\textbf{\ml{ചൂണ്ടുപലക}} — \emph{cūṇṭupalaka} — \textquotedblleft{}a signpost\textquotedblright{}.

\begin{quote}
\ml{ചൂണ്ടുപലകയിൽ} \ml{ഗ്രാമത്തിലേക്കുള്ള} \ml{ദൂരം} \ml{ഉണ്ടായിരുന്നു}.
\end{quote}

\begin{quote}
\emph{cūṇṭupalakayil grāmattilēkkuḷḷa dūraṁ uṇṭāyirunnu.}
\end{quote}

\begin{quote}
The signpost gave the distance to the village.
\end{quote}

Literally a pointing board: \ml{ചൂണ്ടുക}, to point, and \ml{പലക}, a board.

\begin{cousinweb}
\textbf{\ml{ചൂണ്ടുപലക}} (\emph{cūṇṭupalaka}) is a Malayalam compound built from \ml{ചൂണ്ടുക}, to point, which is native Dravidian (DEDR 2658), and \ml{പലക}, a board, from Sanskrit \emph{phalaka}.
\end{cousinweb}

\begin{grammarlens}[title={What you've built}]
One more word for getting around.
\end{grammarlens}

\subsection*{Your turn}
\begin{itemize}
\raggedright
  \item \emph{Say it:} \emph{cūṇṭupalaka}
  \item \emph{Say it:} \emph{cūṇṭupalaka}, once more
  \item \emph{Say it:} \emph{cūṇṭupalakayil grāmattilēkkuḷḷa dūraṁ uṇṭāyirunnu.} — The signpost gave the distance to the village.
  \item \emph{Recall:} say the Malayalam for a footpath, then the Malayalam for the roadside, then say \emph{cūṇṭupalaka} again
\end{itemize}

\begin{tcolorbox}[breakable,colback=teal!4,colframe=teal!35!black,title={Before you move on}]
What does \ml{ചൂണ്ടുപലക} mean? (\textquotedblleft{}A signpost\textquotedblright{}.) Say it once more.
\end{tcolorbox}
\section[\texorpdfstring{\ml{വഴിമാറുക} (vaḻimāṟuka)}{വഴിമാറുക (vaḻimāṟuka)}]{\ml{വഴിമാറുക} (vaḻimāṟuka) — to make way; to step aside}

\label{lesson:ML-C469-valimaruka}



\begin{quote}
Before the new one: say the Malayalam for the roadside, then the Malayalam for a signpost.
\end{quote}

\subsection*{You'll want to know: \ml{വഴിമാറുക}}
\textbf{\ml{വഴിമാറുക}} — \emph{vaḻimāṟuka} — \textquotedblleft{}to make way; to step aside\textquotedblright{}.

\begin{quote}
\ml{ആംബുലൻസ്} \ml{വന്നപ്പോൾ} \ml{എല്ലാ} \ml{വണ്ടികളും} \ml{വഴിമാറി}.
\end{quote}

\begin{quote}
\emph{āmbulansŭ vannappōḷ ellā vaṇṭikaḷuṁ vaḻimāṟi.}
\end{quote}

\begin{quote}
When the ambulance came, all the vehicles made way.
\end{quote}

Past: \ml{വഴിമാറി}, \emph{vaḻimāṟi}.

Literally the way changes: \ml{വഴി} and \ml{മാറുക}, to move aside.

\begin{cousinweb}
\textbf{\ml{വഴിമാറുക}} (\emph{vaḻimāṟuka}) is a native Dravidian compound of \ml{വഴി}, a way (DEDR 5297), and \ml{മാറുക}, to change place (DEDR 4834).
\end{cousinweb}

\begin{grammarlens}[title={What you've built}]
One more word for getting around.
\end{grammarlens}

\subsection*{Your turn}
\begin{itemize}
\raggedright
  \item \emph{Say it:} \emph{vaḻimāṟuka}
  \item \emph{Say it:} \emph{vaḻimāṟuka}, once more
  \item \emph{Say it:} \emph{āmbulansŭ vannappōḷ ellā vaṇṭikaḷuṁ vaḻimāṟi.} — When the ambulance came, all the vehicles made way.
  \item \emph{Recall:} say the Malayalam for the roadside, then the Malayalam for a signpost, then say \emph{vaḻimāṟuka} again
\end{itemize}

\begin{tcolorbox}[breakable,colback=teal!4,colframe=teal!35!black,title={Before you move on}]
What does \ml{വഴിമാറുക} mean? (\textquotedblleft{}To make way; to step aside\textquotedblright{}.) Say it once more.
\end{tcolorbox}
\section[\texorpdfstring{\ml{മുറിച്ചുകടക്കുക} (muṟiccukaṭakkuka)}{മുറിച്ചുകടക്കുക (muṟiccukaṭakkuka)}]{\ml{മുറിച്ചുകടക്കുക} (muṟiccukaṭakkuka) — to cross over (a road)}

\label{lesson:ML-C469-muriccukatakkuka}



\begin{quote}
Before the new one: say the Malayalam for a signpost, then the Malayalam for to make way.
\end{quote}

\subsection*{You'll want to know: \ml{മുറിച്ചുകടക്കുക}}
\textbf{\ml{മുറിച്ചുകടക്കുക}} — \emph{muṟiccukaṭakkuka} — \textquotedblleft{}to cross over (a road)\textquotedblright{}.

\begin{quote}
\ml{മുത്തശ്ശി} \ml{പതുക്കെ} \ml{വഴി} \ml{മുറിച്ചുകടന്നു}.
\end{quote}

\begin{quote}
\emph{muttaśśi patukke vaḻi muṟiccukaṭannu.}
\end{quote}

\begin{quote}
Grandmother slowly crossed the road.
\end{quote}

Past: \ml{മുറിച്ചുകടന്നു}, \emph{muṟiccukaṭannu}.

Two verbs: \ml{മുറിച്ച്}, having cut, and \ml{കടക്കുക}, to cross. You cut across the road.

\begin{cousinweb}
\textbf{\ml{മുറിച്ചുകടക്കുക}} (\emph{muṟiccukaṭakkuka}) is a native Dravidian compound of \ml{മുറിക്കുക}, to cut (DEDR 5008), and \ml{കടക്കുക}, to cross (DEDR 1109).
\end{cousinweb}

\begin{grammarlens}[title={What you've built}]
One more word for getting around.
\end{grammarlens}

\subsection*{Your turn}
\begin{itemize}
\raggedright
  \item \emph{Say it:} \emph{muṟiccukaṭakkuka}
  \item \emph{Say it:} \emph{muṟiccukaṭakkuka}, once more
  \item \emph{Say it:} \emph{muttaśśi patukke vaḻi muṟiccukaṭannu.} — Grandmother slowly crossed the road.
  \item \emph{Recall:} say the Malayalam for a signpost, then the Malayalam for to make way, then say \emph{muṟiccukaṭakkuka} again
\end{itemize}

\begin{tcolorbox}[breakable,colback=teal!4,colframe=teal!35!black,title={Before you move on}]
What does \ml{മുറിച്ചുകടക്കുക} mean? (\textquotedblleft{}To cross over (a road)\textquotedblright{}.) Say it once more.
\end{tcolorbox}
